\documentclass[11pt]{article}
\usepackage[margin=1in]{geometry}
\usepackage{amssymb}
\usepackage{hyperref}
\title{TLA+ Study Group — Week 04}
\author{Bartosz}
\date{\today}
\begin{document}\maketitle

\section*{Brief Description of the Topics}

\begin{itemize}
  \item { Description of PAXOS a fault-tolerant distributed transaction-commit algorithm }
\end{itemize}

New syntax:
\begin{description}
  \item[CHOOSE] { 
    \begin{verbatim}
      CHOOSE v \in S : P equals
    \end{verbatim}
    if S is a singleton set then v is the value of the element of the singleton
    set; otherwise v is unspecified but P holds. If P cannot hold TLA will mark
    an error. One should only use it in the case of singletons or when, in a
    larger expression, the value doesn't depend on which $v$ was chosen
  }
  \item[LET] { 
    \begin{verbatim}
    LET value == $expr
    IN $expr
    \end{verbatim}
  }
  \item[set comprehension $v \in S : P$] {
    \begin{verbatim} {
      { v \in S : P }
    \end{verbatim}
    subset of S consisting of all $v$ satisfying ,P for example:
    \begin{verbatim}
      { n \in Nat : n > 17 } (* = { 18, 19, 20 ... }
    \end{verbatim}
  }
  \item[set comprehension $e : v \in S$] {
    \begin{verbatim} {
      { e : v \in S }
    \end{verbatim}
    for example
    \begin{verbatim}
    {n^2 : n \in Nat} (* { 0, 1, 4, 9... }
    \end{verbatim}
  }
  \item[generalization of EXECPT] {
    \begin{verbatim}
    aState' = [aState EXCEPT ![m.ins][acc].mbal = m.bal,
                             ![m.ins][acc].bal  = m.bal,
                             ![m.ins][acc].val  = m.val
    \end{verbatim}
    is equivalent to
    \begin{verbatim}
    aState[m.ins][acc].mbal = m.bal;
    aState[m.ins][acc].bal = m.bal;
    aState[m.ins][acc].val = m.val;
    \end{verbatim}

    i.e a "copy method" for records that allows for arbitrary nesting and
    chaining
  }
\end{description}

\section*{What stood out}


\section*{Open Questions}

\end{document}
